\documentclass[11pt]{article}
\usepackage{mystyle}

\title{Rosen --- \S 6.5: Generalized Permutations and Combinations\\ \large Selected Exercises}
\author{Alston}
\date{Semester 2, 2026}

\begin{document}
\maketitle

% ======================================================================
\begin{exercise}{30}
    How many different strings can be made from the letters in
    MISSISSIPPI, using all the letters?
\end{exercise}

% ======================================================================
\begin{exercise}{40}
    How many ways are there to travel in $xyzw$ space from the origin
    $(0,0,0,0)$ to the point $(4,3,5,4)$ by taking steps one unit in
    the positive $x$, positive $y$, positive $z$, or positive $w$
    direction?
\end{exercise}

% ======================================================================
\begin{exercise}{46}
    A shelf holds 12 books in a row. How many ways are there to
    choose five books so that no two adjacent books are chosen?
    [\emph{Hint}: Represent the books that are chosen by bars and the
    books not chosen by stars. Count the number of sequences of five
    bars and seven stars so that no two bars are adjacent.]
\end{exercise}

% ======================================================================
\begin{exercise}{48}
    \textit{Recall: Theorem 3 states that the number of different
    permutations of $n$ objects, where there are $n_1$
    indistinguishable objects of type 1, $n_2$ indistinguishable
    objects of type 2, \dots, and $n_k$ indistinguishable objects of
    type $k$, is $n!/(n_1! n_2! \cdots n_k!)$. Theorem 4 states that
    the number of ways to distribute $n$ distinguishable objects into
    $k$ distinguishable boxes so that $n_i$ objects are placed into
    box $i$, $i = 1, 2, \dots, k$, equals
    $n!/(n_1! n_2! \cdots n_k!)$.}

    Prove Theorem 4 by first setting up a one-to-one correspondence
    between permutations of $n$ objects with $n_i$ indistinguishable
    objects of type $i$, $i = 1, 2, 3, \dots, k$, and the
    distributions of $n$ objects in $k$ boxes such that $n_i$ objects
    are placed in box $i$, $i = 1, 2, 3, \dots, k$, and then applying
    Theorem 3.
\end{exercise}

% ======================================================================
\begin{exercise}{49}
    \textit{Recall: Theorem 2 states that there are
    $\binom{n+r-1}{r} = \binom{n+r-1}{n-1}$ $r$-combinations from a
    set with $n$ elements when repetition of elements is allowed.}

    In this exercise we will prove Theorem 2 by setting up a
    one-to-one correspondence between the set of $r$-combinations
    with repetition allowed of $S = \{1, 2, 3, \dots, n\}$ and the
    set of $r$-combinations of the set $T = \{1, 2, 3, \dots,
    n+r-1\}$.
    \begin{enumerate}[label=(\alph*)]
        \item Arrange the elements in an $r$-combination, with
        repetition allowed, of $S$ into an increasing sequence
        $x_1 \leq x_2 \leq \cdots \leq x_r$. Show that the sequence
        formed by adding $k-1$ to the $k$th term is strictly
        increasing. Conclude that this sequence is made up of $r$
        distinct elements from $T$.
        \item Show that the procedure described in (a) defines a
        one-to-one correspondence between the set of
        $r$-combinations, with repetition allowed, of $S$ and the
        $r$-combinations of $T$.
        [\emph{Hint}: Show the correspondence can be reversed by
        associating to the $r$-combination $\{x_1, x_2, \dots, x_r\}$
        of $T$, with $1 \leq x_1 < x_2 < \cdots < x_r \leq n+r-1$,
        the $r$-combination with repetition allowed from $S$, formed
        by subtracting $k-1$ from the $k$th element.]
        \item Conclude that there are $\binom{n+r-1}{r}$
        $r$-combinations with repetition allowed from a set with $n$
        elements.
    \end{enumerate}
\end{exercise}

% ======================================================================
\begin{exercise}{59}
    How many ways are there to distribute five balls into three boxes
    if each box must have at least one ball in it if
    \begin{enumerate}[label=(\alph*)]
        \item both the balls and boxes are labeled?
        \item the balls are labeled, but the boxes are unlabeled?
        \item the balls are unlabeled, but the boxes are labeled?
        \item both the balls and boxes are unlabeled?
    \end{enumerate}
\end{exercise}

% ======================================================================
\begin{exercise}{61}
    Suppose that a weapons inspector must inspect each of five
    different sites twice, visiting one site per day. The inspector
    is free to select the order in which to visit these sites, but
    cannot visit site $X$, the most suspicious site, on two
    consecutive days. In how many different orders can the inspector
    visit these sites?
\end{exercise}

% ======================================================================
\begin{exercise}{63}
    Prove the Multinomial Theorem: If $n$ is a positive integer, then
    \[
    (x_1 + x_2 + \cdots + x_m)^n = \sum_{n_1+n_2+\cdots+n_m=n}
    C(n; n_1, n_2, \dots, n_m)\, x_1^{n_1} x_2^{n_2} \cdots
    x_m^{n_m},
    \]
    where
    \[
    C(n; n_1, n_2, \dots, n_m) = \frac{n!}{n_1!\, n_2! \cdots n_m!}
    \]
    is a multinomial coefficient.
\end{exercise}

\end{document}